\section{Development Tooling and AI Assistance}
\label{sec:tools-and-technologies}

The rationale for the platform's technology stack is architectural rather than methodological, and is therefore presented with the other design decisions in \cref{sec:technology-selection}.
This section records two aspects of how the artifact was produced: the tooling that kept it continuously validated during development, and the use of generative-AI assistants.

Throughout development the codebase was kept continuously buildable and testable through two independent continuous-integration pipelines, one for each sub-project, so that every change was validated automatically.
The engineering practices behind this, including the automated test suite and the Windows build environment required by the eye-tracking hardware (C1), are detailed in \cref{ch:implementation}.

In keeping with contemporary software practice, generative-AI assistants were used in a limited, supporting role during the project.
Their use was bounded to three tasks: coding assistance for the software artifact, feedback and rephrasing on prose the authors had already written, and a sounding board for discussing design trade-offs the authors had already formed.
All AI-assisted output was reviewed, verified, and (where it was code) tested by the authors, who remain responsible for the result; the analysis, the design decisions, and the prose of this thesis are the authors' own.
A full declaration of generative-AI use, including the specific tools and the tasks for which they were used, is provided in \cref{app:ai-declaration}.
